\begin{recipe}
[% 
    preparationtime = {\unit[25]{min}},
    bakingtime = {\unit[50]{min}},
    bakingtemperature = {\protect\bakingtemperature{
        topbottomheat=\unit[160]{\textcelcius}
    }},
    portion = {\portion{8}},
    source = {\protect\recipesource{https://www.instagram.com/p/DTqvDXFDOKK/}{Instagram: raikas\_show}},
    calory = {756kcal | 7g Protein | 103g KH | 38g Fett},
    price = {1,35€},
    %rating = {1},
]
{Sticky Toffee Pudding}

%\graph{% pictures
%    big=pic/dessert/08_sticky-toffee-pudding
%}

\introduction{%
Saftig, karamellig und mit einer intensiven Toffee-Sauce überzogen – Sticky Toffee Pudding ist eines dieser Desserts, das man eigentlich nicht teilen will. Warm serviert absolut gefährlich.
}

\ingredients{
    \textbf{Teig} & \\
    125 g & Butter\index{Milchprodukte \& Alternativen!Butter}\\
    250 g & Datteln\index{Obst!Dattel}\\
    250 ml & Heißes Wasser\\
    150 g & Brauner Zucker\\
    70 g & Weißer Zucker\\
    90 g & Melasse\\
    200 g & Saure Sahne\index{Milchprodukte \& Alternativen!Saure Sahne}\\
    275 g & Weizenmehl\index{Getreide!Weizenmehl}\\
    2 & Eier\index{Eier!Ei}\\
    4 g & Natron\\
    & Vanilleextrakt, Zimt, 1 TL Salz\\
    \\
    \textbf{Sauce} & \\
    125 g & Butter\\
    150 g & Brauner Zucker\\
    60 g & Melasse\\
    250 ml & Sahne\index{Milchprodukte \& Alternativen!Sahne}\\
    & Vanilleextrakt, Prise Salz
}

\preparation{%
    \step%
    Backofen auf 160 °C Ober-/Unterhitze vorheizen.\\
    
    \step%
    Datteln klein hacken und mit heißem Wasser übergießen. 10 Minuten einweichen lassen.
    
    \step%
    Butter mit braunem und weißem Zucker cremig schlagen. Eier unterrühren, dann Saure Sahne, Melasse und Gewürze zugeben.
    
    \step%
    Mehl und Natron mischen und unterheben. Zum Schluss die eingeweichten Datteln samt Flüssigkeit unterrühren.
    
    \step%
    Teig in eine gefettete Form geben und 45--60 Minuten backen (Stäbchenprobe).
    
    \step%
    Für die Sauce Butter, Zucker und Melasse erhitzen. Sahne zugeben und 3--5 Minuten köcheln lassen, bis die Sauce leicht eindickt. Mit Vanille und Salz abschmecken.
    
    \step%
    Warmen Kuchen großzügig mit Sauce übergießen und servieren.
}

\suggestion[Servieren]{%
    Am besten warm mit Vanilleeis oder etwas geschlagener Sahne servieren.
}

\end{recipe}